\documentclass[11pt]{article}
\usepackage{amsmath,amssymb,amsthm}
\usepackage[margin=1in]{geometry}
\usepackage{hyperref}

\newtheorem{theorem}{Theorem}[section]
\newtheorem{proposition}[theorem]{Proposition}
\newtheorem{lemma}[theorem]{Lemma}
\newtheorem{corollary}[theorem]{Corollary}
\theoremstyle{definition}
\newtheorem{definition}[theorem]{Definition}
\newtheorem{example}[theorem]{Example}
\theoremstyle{remark}
\newtheorem{remark}[theorem]{Remark}

\newcommand{\Stab}{\mathrm{Stab}}
\newcommand{\LO}{\mathbf{LO}}
\newcommand{\Fin}{\mathrm{Fin}}
\newcommand{\Sub}{\mathrm{Sub}}
\newcommand{\Dec}{\mathrm{Dec}}
\newcommand{\WF}{\mathrm{WF}}
\newcommand{\Acc}{\mathrm{Acc}}
\newcommand{\Hom}{\mathrm{Hom}}

\title{Dilators and ptykes as categories of embeddings,\\ constructively}
\author{Notes on a formalisation (draft, 19 November 2026)}
\date{}

\begin{document}
\maketitle

\begin{abstract}
We describe a constructive development, checked mechanically, of Girard's
dilators and ptykes~\cite{Gir81,GirPTLC} following the approach of
\emph{categories of embeddings}~\cite{Coq89}. The central result is that the category
$\Stab(C)$ of stable functors on a category of embeddings $C$ with good
finite objects is again such a category; iterating from the category of
linear orders gives dilators, ptykes of type~2, of type~3, and so on.
Well-foundedness is inductive accessibility, and the main development uses
no classical principle. We give examples (polynomial dilators, $2^X$, $\varepsilon_X$, the
evaluation ptykes, and the Bachmann--Howard collapse $D\mapsto\vartheta D$
as a type-2 ptyx), the constructive halves of Girard's $\Pi^1_2$ and
$\Pi^1_3$ completeness theorems (and, in a separate classical segment, the
full $\Pi^1_2$ theorem), and a candidate type-3 collapse aimed at
$\Pi^1_2$-comprehension. With univalence, finite stable functors up to
isomorphism form a small set with decidable equality, and dilators and
ptykes of type~2 form sets.
\end{abstract}

\section{Setting}

We work in a constructive dependent type theory (with a hierarchy of
universes, function extensionality and unique choice; propositional
truncation is used for mere existence).  All
orders are \emph{decidable}: a linear order is a type $A$ with a
comparison $A\times A\to\{<,=,>\}$ satisfying reflexivity, flip,
transitivity and $x=y$ whenever the comparison says so.  A morphism of
linear orders is a strictly increasing map (an embedding).  An order is
\emph{well-founded} if every element is accessible (inductively defined
$\Acc$).  The main development uses no excluded middle; in several places
decidability that is usually assumed is \emph{derived}. Two extensions are
kept separate, as parameters of the segments that use them: univalence and
replacement (Section~\ref{sec:ua}), and excluded middle with dependent
choice, used only to complete Girard's $\Pi^1_2$ theorem
(Section~\ref{sec:girard}).

\section{Categories of embeddings}

In~\cite{Coq89} a stable category of embeddings is the ind-completion of a
\emph{category of events}, a categorical version of Winskel's event
structures~\cite{Win}; the stable functors on it are those of Girard's
$\Pi^1_2$-logic~\cite{Gir81}. We use the following constructive variant.

\begin{definition}[Category of embeddings]
A \emph{category of embeddings} $C$ consists of a category whose morphisms
are monomorphisms, with pull-backs, together with a small type $\Fin$ of
codes and a map $\mathrm{fin}:\Fin\to C$ (the finite objects, not
necessarily injectively coded), such that any two maps from finite objects
into an object merely factor through a common finite object
(\emph{directedness}).
\end{definition}

\begin{definition}[Good finite objects]
A category of embeddings $C$ has \emph{good finite objects} (a CEF) when:
\begin{enumerate}
\item hom-sets between finite objects are Bishop-finite, with decidable
  equality;
\item each finite object has a finite list of subobjects covering, up to
  image-equivalence, every map from a finite object into it;
\item every listed subobject $s:b\to a$ has a \emph{self-amalgamation}: two
  maps $w,w':a\to d$ agreeing on $s$ whose pull-back factors through $s$;
\item finite objects are rigid (their only endomorphism is the identity);
\item finite \emph{multicoproducts} in the sense of Diers~\cite{Die80}
  (see also~\cite{Osm1,Osm2}): for finite
  objects $a,b$ a finite (junk-tolerant) list of cospans $a\to c\leftarrow
  b$ such that every pair of maps $a\to d\leftarrow b$ into a finite object
  factors through exactly one listed cospan, by exactly one map.
\end{enumerate}
\end{definition}

\begin{example}
Finite ordinals with strictly increasing maps (coded by words over
$\{0,1\}$) form a CEF: pull-backs are intersections of images, amalgamation
doubles the points outside a subobject, the multicoproducts of $[a]$ and
$[b]$ are the jointly surjective shuffles.  The category $\LO$ of all
decidable linear orders with finite objects $[n]$ is a CEF, the finiteness
data being transported from the finite ordinals.
\end{example}

\begin{definition}[Stable functors]
A \emph{stable functor} on $C$ is a functor $X:C\to\LO$ that is
\emph{continuous} (every point of $X(A)$ comes from some finite object
mapping to $A$) and \emph{stable} (it preserves pull-backs, in image
form). Morphisms of stable functors are natural transformations that are
\emph{cartesian} (naturality squares are pull-backs); they are represented
by their restriction to finite objects, so that hom-sets are small, and
extended to all objects by unique choice.
\end{definition}

In Diers's terms~\cite{Osm1}, stability says that every point factors
uniquely through a candidate; below, the candidate is the least support.

\begin{proposition}[Decidable supports]
Let $s:B\to A$ be a map with a self-amalgamation $w,w':A\to D$. For every
stable functor $X$ and $x\in X(A)$,
\[ x\in \mathrm{Im}\,X(s)\iff X(w)x=X(w')x ,\]
an equation in the linear order $X(D)$; hence membership in the image of
$X(s)$ is decidable.
\end{proposition}

This replaces the classical assumption of decidable supports (or decidable
degeneracies): it is a consequence of stability.

\begin{theorem}[Normal forms]
Let $C$ be a CEF and $X$ a stable functor. Every point $x\in X(\mathrm{fin}\,c)$
has a least support, found by a finite search over the listed subobjects,
and normal forms are unique: if $X(k)x_0=X(k')x_0$ for a point $x_0$ at its
least support then $k=k'$.
\end{theorem}

\begin{definition}[Codes]
A \emph{code} consists of $m$ generators with finite objects
$\mathrm{gen}(i)$ and, at each finite level $d$, a linear order on the
normal forms $(i,k:\mathrm{gen}(i)\to d)$, invariant under maps and
satisfying the order laws. A code $K$ \emph{decodes} to a stable functor
$\mathrm{Decode}\,K$ whose points at $A$ are pairs
$(i,k:\mathrm{fin}(\mathrm{gen}\,i)\to A)$, without any quotient.
\end{definition}

\begin{theorem}[Covering]
Every point of a stable functor $X$ is the image, under a cartesian
morphism $\mathrm{Decode}\,K\to X$, of a point of a decoded code $K$.
\end{theorem}

\begin{theorem}[Main theorem]
If $C$ is a CEF then $\Stab(C)$, the category of stable functors on $C$
and cartesian morphisms, with codes as finite objects, is again a CEF.
\end{theorem}

The proof assembles: pull-backs computed pointwise; directedness by
merging the generators of two codes, removing duplicates (decidable);
Bishop-finiteness of morphisms between decoded codes (a morphism is
determined by the images of the generators, and validity of an assignment is
decidable using the multicoproducts of $C$); subobjects as subsets of
generators; amalgamation by a lexicographic product with $\{L,R\}$;
rigidity from a constructive pigeonhole principle (an order-preserving
self-map of a Bishop-finite linear order is the identity); and the
multicoproducts of two decoded codes $K_1,K_2$: a listed cospan is given by
a subset of the generators of $K_2$ to be kept, a \emph{table} of
comparisons between generators of $K_1$ and kept generators of $K_2$ at the
listed joint types of their finite objects, which is valid (decidably), and
the resulting sum code.

\begin{corollary}[The hierarchy]
$\LO$, $\Stab(\LO)$ (pre-dilators), $\Stab(\Stab(\LO))$ (pre-ptykes of
type 2), $\Stab^3(\LO)$, \dots{} are all categories of embeddings with good
finite objects.
\end{corollary}

\section{Dilators and ptykes}

\begin{definition}\label{def:dil}
A \emph{pre-dilator} is a stable functor $D$ on $\LO$. It is
\emph{monotone} if for every point $x\in D([n])$ and embeddings
$k,k':[n]\to A$ with $k(j)\le k'(j)$ for all $j<n$,
\[ D(k)\,x\ \le\ D(k')\,x \quad\text{in } D(A). \]
A pre-dilator $D$ is a \emph{dilator} if $D(A)$ is well-founded whenever
$A$ is. A pre-ptyx of type~2 is a stable functor $F$ on $\Stab(\LO)$; it is
a \emph{ptyx} if $F(D)$ is well-founded whenever $D$ is a dilator. Similarly
at higher types.
\end{definition}

Girard's predilators~\cite{Gir81,GirPTLC,GN} are the \emph{monotone}
pre-dilators, and his dilators the monotone pre-dilators preserving
well-orders: monotonicity is part of his definition. Girard shows that it
follows from the preservation of well-orders; this holds constructively, so
the two notions of dilator agree.

\begin{theorem}[Dilators are monotone]\label{thm:mono}
Every pre-dilator preserving well-orders is monotone.
\end{theorem}

\begin{proof}[Sketch]
The comparison of $D(k)x$ and $D(k')x$ is decidable, so it suffices to
refute $D(k)x>D(k')x$. By directedness we may assume $A=[c]$ finite, and,
changing one coordinate at a time from the top, that $k$ and $k'$ differ
only at one coordinate, $a=k(r)<b=k'(r)$. In the well-order
$B=\omega\times\omega\times\omega$ (lexicographic) the embeddings
$\psi_i:[c]\to B$ given by $v\mapsto(v,0,v)$ for $v<a$ or $v>b$,
$v\mapsto(a,i,v)$ for $a\le v<b$, and $b\mapsto(a,i+1,a)$ satisfy
$\psi_i\circ k'=\psi_{i+1}\circ k$. Hence $D(\psi_i\circ k)x$ would be an
infinite descending sequence in the well-order $D(B)$.
\end{proof}

For a pre-dilator given by a code, monotonicity on the generators (if
$k\le k'$ pointwise then $(i,k)\le(i,k')$) implies monotonicity on all
points, since every point is the image of a generator.

\begin{example}[First examples]
The identity $X\mapsto X$ is a monotone dilator; $X\mapsto X^{\mathrm{op}}$
is a pre-dilator that is neither monotone ($k<k'$ is reversed) nor a dilator
($\omega^{\mathrm{op}}$ has no accessible point). Evaluation
$D\mapsto D(Z)$ at a well-founded $Z$ is a type-2 ptyx, e.g.\ at $\omega$.
\end{example}

\begin{example}[Polynomial dilators and $2^X$]
Every polynomial built from the identity, constants, sums, products and
binary exponentiation $2^D$ is a dilator on all linear orders; in
particular $X\mapsto 2^X$ (Girard's well-ordering principle, classically
equivalent to $\mathrm{ACA}_0$~\cite{GirPTLC}), towers $2^{2^{\cdot^{\cdot^X}}}$, and the
evaluation $D\mapsto D(2^\omega)$ is a type-2 ptyx.
\end{example}

\begin{example}[$\varepsilon_X$]
Cantor normal forms over the symbols $\varepsilon_x$ ($x\in X$) form a
decidable linear order; $X\mapsto\varepsilon_X$ is a stable functor, and a
dilator: if $X$ is well-founded then so is $\varepsilon_X$. Classically
this principle is equivalent to $\mathrm{ACA}_0^+$ (Afshari--Rathjen~\cite{AR},
Marcone--Montalb\'an~\cite{MM}; see also the survey~\cite{Rat21}).
\end{example}

\begin{theorem}[Bachmann--Howard collapse]
For every pre-dilator $D$, Freund's order $\vartheta D$~\cite{Fre19} (terms
$\vartheta(\sigma)$ with $\sigma\in D(\vartheta D)$, compared by
\[\vartheta\sigma<\vartheta\tau\iff \vartheta\sigma\le\text{some support
element of }\tau,\ \text{or}\ \sigma<\tau\text{ and all support elements of
}\sigma\text{ are }<\vartheta\tau)\]
is a decidable linear order; $D\mapsto\vartheta D$ is a stable functor on
$\Stab(\LO)$; and it is a type-2 ptyx: $\vartheta D$ is well-founded for
every dilator $D$.
\end{theorem}

Classically, ``$\vartheta D$ is well-founded for every dilator $D$'' is
equivalent to $\Pi^1_1$-comprehension (Freund~\cite{Fre19}; see
also~\cite{AFW} for Girard's $\Lambda$ and P\"appinghaus's functor~$J$). The constructive proof
takes the accessible part $W$ of $\vartheta D$, uses that $D(W)$ is
well-founded, and argues by induction on the image of $\sigma$ in $D(W)$,
with a semantic lemma identifying the term comparison with the comparison
in $D$ of the image of a well-behaved sub-order.

\section{Girard's completeness theorems}\label{sec:girard}

Both theorems are due to Girard~\cite{Gir81,GirPTLC}. We formalise
constructively the directions that need no classical principle, and, for
$\Pi^1_2$, the converse in a classical segment.

\begin{theorem}[$\Pi^1_2$]
To a tree $T$ (a $\Pi^1_2$ statement $\forall f\exists g\forall n\,R$)
one associates a pre-dilator $D_T$ (labelled trees ordered by the
Kleene--Brouwer order) such that $D_T$ is a dilator iff all its attempt
trees are well-founded, and a rank of some $U_f$ into a well-founded order
refutes that $D_T$ is a dilator.
\end{theorem}

\begin{theorem}[$\Pi^1_3$]
The same construction one type up gives, from a $\Pi^1_3$ statement, a
type-2 pre-ptyx (trees whose nodes carry partial natural embeddings into
the dilator argument, compared on the listed joint types), which is a ptyx
iff the corresponding trees are well-founded for every dilator.
\end{theorem}

\begin{theorem}[$\Pi^1_4$]
Fix a sequence of codes of finite dilators, each decoding to a dilator. A
$\Pi^1_4$ statement $\forall f\,\exists g\,\forall h\,\exists k\ldots$ gives a
type-3 pre-ptyx $G$: its value at a type-2 pre-ptyx $F$ is the
Kleene--Brouwer order of finite attempts to embed, naturally, the finite
part of the type-2 pre-ptyx $F_u$ of the previous theorem into $F$
(comparisons checked on the listed multicoproducts of codes). Then $G$ is a
type-3 ptyx iff every 2-ptyx has a well-founded attempt tree; and if $F_f$
is a 2-ptyx for some $f$, then $G$ is not a type-3 ptyx.
\end{theorem}

The converse directions use bar induction. For $\Pi^1_2$ they are
formalised in a separate segment taking excluded middle (for
propositions) and dependent choice as parameters:

\begin{theorem}[Girard's $\Pi^1_2$ theorem, classically]
Assume excluded middle and dependent choice, let the trees $U_f$ be
prefix-closed and let the enumeration of finite sequences list every $t$ at
an index $j$ with $|t|\le j+1$. Then $D_T$ is a dilator iff every $U_f$ has
an infinite branch.
\end{theorem}

The only classical ingredient is bar induction: a prefix-closed tree over a
well-founded label order with no infinite branch has a well-founded
Kleene--Brouwer order. Kleene--Brouwer accessibility of the subtree at a
node is a proposition, closed under ``all children'' (constructively);
if the root fails, excluded middle gives a failing child at each step and
dependent choice an infinite branch. Then: if some $U_f$ has no infinite
branch, its own Kleene--Brouwer order is a rank into a well-order, which
refutes that $D_T$ is a dilator; conversely an infinite branch of an
attempt tree over a well-order $A$ gives $f$ and a rank of $U_f$ into
$1+A$, which refutes the infinite branches of $U_f$. With excluded middle,
propositions are small (impredicativity), so this segment is a
verification of Girard's classical theorem on top of the constructive
development, not part of it. The $\Pi^1_3$ and $\Pi^1_4$ converses (which
need Girard's classical type-2 completeness theorem) are not formalised.
A positive instance (the tree of strictly decreasing sequences) is proved
constructively to give a dilator.

The hypothesis that the listed codes decode to dilators is discharged
for \emph{monotone} codes (Definition~\ref{def:dil}; for a code it suffices
to check the generators). Not every code is a dilator:
$X\mapsto X^{\mathrm{op}}$ is a one-generator code which is not. By
Theorem~\ref{thm:mono} and the next theorem, a code decodes to a dilator iff
it is monotone, which is decidable.

\begin{theorem}[Monotone finite pre-dilators are dilators]
If a code $K$ is monotone then $\mathrm{Decode}\,K$ is a dilator.
\end{theorem}

The proof uses almost-full relations: following Coquand's direct proof of
the intuitionistic Ramsey theorem~\cite{CoqRamsey}, if a tree $p$ secures
$A\vee R$ and $q$ secures $B\vee S$ for $n$-ary $R,S$, then $p\otimes_n q$
secures $A\vee B\vee(R\wedge S)$; hence almost-full relations are closed
under intersection. On a well-founded linear order, $\le$ is almost full;
almost-fullness pulls back along maps; so ``same generator and pointwise
$\le$'' is almost full on $\mathrm{Decode}\,K(A)$; by monotonicity it is
contained in $\le$, and an almost-full relation contained in $\le$ makes a
linear order well-founded.

\begin{corollary}[A statement beyond the type theory]
Take for $S$ the $\Pi^1_4$ axiom of $\Pi^1_2$-$\mathrm{CA}_0$ and let $G_S$ be
the type-3 pre-ptyx of the previous theorem. The formalised implication
``$G_S$ is a type-3 ptyx $\Rightarrow$ no $F_f$ is a 2-ptyx'', followed by
the classical type-2 completeness theorem and Kleene's normal form, gives
\[ G_S\ \text{is a type-3 ptyx}\ \Longrightarrow\ \Pi^1_2\text{-}\mathrm{CA}_0 .\]
The statement ``$G_S$ is a type-3 ptyx'' is expressible in the type theory
used here. Since this type theory is predicative, of strength well below
$\Pi^1_2$-$\mathrm{CA}_0$, the statement is not provable in it (a
metatheorem, via G\"odel's second incompleteness theorem; not formalised).
\end{corollary}

\section{With univalence: finite objects up to isomorphism}\label{sec:ua}

In this section we assume in addition the univalence axiom and Rijke's
replacement principle~\cite{Rij} (the image of a map from a small type
into a locally small type is small), as in the treatment of Scott domains
in the same library. The point is that \emph{rigidity} of finite objects
interacts well with univalence: the finite objects up to isomorphism form a
small \emph{set}, obtained without choosing codes, in the same way as the
compact elements of a Scott domain form a small basis.

\begin{theorem}[Univalence of $\Stab(C)$]
For every category of embeddings $C$ and stable functors $X,Y$, the type of
paths $X=Y$ is equivalent to the type of natural isomorphisms $X\cong Y$
(isomorphisms of linear orders at every object, natural in $C$), and also
to the type of natural isomorphisms between the restrictions of $X$ and $Y$
to the finite objects. The latter type is small, so $\Stab(C)$ is locally
small.
\end{theorem}

The first equivalence is the structure identity principle for the family
of linear orders together with its action on maps: naturality forces the
transported action to coincide with the given one, and the remaining
components are propositions. For the second, a natural isomorphism on
finite objects and its inverse extend to all objects (continuity), and the
extensions are inverse by uniqueness of extensions. At level~$0$ the same
argument shows that the finite linear orders form a small type equivalent
to $\mathbb{N}$: since $[n]$ is rigid, an isomorphism $A\cong[n]$ is unique
when it exists, so ``$A$ is finite'' needs no truncation.

\begin{theorem}[Finite stable functors up to isomorphism]
Let $C$ be a CEF and let $\mathrm{FinDil}_C$ be the image of
$\mathrm{Decode}:\mathrm{Codes}\to\Stab(C)$, small by replacement. Then
\begin{enumerate}
\item $\mathrm{FinDil}_C$ is a set: since a decoded code is rigid, the
  natural isomorphisms out of it form a proposition, hence so do its
  paths;
\item $\mathrm{FinDil}_C$ has decidable equality: by rigidity \emph{any}
  pair of maps $\mathrm{Decode}\,K\rightleftarrows\mathrm{Decode}\,K'$ is
  an isomorphism (both composites are endomorphisms, hence identities),
  and the existence of such maps is decided by two searches in the finite
  hom-sets.
\end{enumerate}
\end{theorem}

\paragraph{Dilators form a set.} The same mechanism applies beyond finite
objects, exactly as for the ordinals (the type of ordinals is a set because
a well-order has no non-trivial automorphism~\cite[\S10.3]{HoTT}). If $f,g$
are mutually inverse order-preserving maps of a well-founded linear order,
then $f=\mathrm{id}$, by accessibility induction: if $f(x)<x$ then
$f(f(x))=f(x)$ by induction, so $f(x)=x$; if $f(x)>x$ then $g(x)<x$, and
$f(g(x))=g(x)$ gives $g(x)=x$. For a dilator $X$, every $X([n])$ is
well-founded, so two natural isomorphisms out of $X$ agree on the finite
objects, hence everywhere.

\begin{theorem}
For a dilator $X$ and any stable functor $Y$ on $\LO$, the type of paths
$X=Y$ is a proposition. Hence the type of dilators,
$\sum_{X:\Stab(\LO)}\mathrm{IsDil}(X)$, is a set.
\end{theorem}

The type of pre-dilators is not a set: the constant functor with value
$\mathbb{Z}$ has the shift as a non-trivial natural automorphism. Thus
being a dilator, a well-foundedness condition, is what removes the higher
structure. Note that the rigidity of finite linear orders is stronger: every
order-preserving \emph{self-map} is the identity, which is what makes
equality of finite stable functors decidable. At type~$2$ the argument does
not apply directly, since a ptyx is evaluated also at finite pre-dilators
that are not dilators.

More precisely, for any category of embeddings $C$ and any stable functor
$X$ whose values at the finite objects are well-founded, paths $X=Y$ are
equivalent to isomorphisms $X\cong Y$ and both are propositions: such
functors form a univalent \emph{gaunt} category (its only isomorphisms are
identities). It is not a preorder ($\mathrm{Id}\to\mathrm{Id}+\mathrm{Id}$
has two embeddings).

\paragraph{Ptykes of type 2.} To go one type up, the dilators must
themselves form a category of embeddings, whose finite objects form a
\emph{small} type. Taken literally, ``$\mathrm{Decode}\,K$ is a dilator''
quantifies over all linear orders and is a large proposition. But by
Theorem~\ref{thm:mono} it is equivalent to monotonicity, which for a
finite code is small and decidable (it is checked at the joint types of
each generator with itself). So the finite objects are the monotone codes
and the objects the dilators (which are monotone).
Monotonicity and well-foundedness pass to sub-functors, which gives
pull-backs and directedness. A ptyx of type~2 in Girard's sense is a stable
functor on this category which is well-founded at every object, and:

\begin{theorem}
Dilators and ptykes of type 2 (in Girard's sense) form univalent gaunt
categories; the type of dilators and the type of ptykes of type 2 are sets.
\end{theorem}

\begin{remark}
Here monotonicity does real work: it turns the large proposition ``this
finite pre-dilator is a dilator'' into a small decidable one.
\end{remark}

Going further up requires a small criterion for the finite objects of type
2 (an analogue of monotonicity), which we do not have.

\paragraph{Enumeration.} A code has a comparison at \emph{every} finite
level, so codes are not finite data. They are, up to isomorphism: by the
multicoproducts, a code is determined by a finite \emph{table} giving, for
two generators $i,i'$ and a listed joint type $t$ of their finite objects,
the comparison of the two normal forms at $t$. Two points at any level are
compared at the least listed joint type through which they factor (found
by search, and invariant under maps by uniqueness of factorisations).
Validity of a table reduces to finitely many checks: reflexivity at the
generators, flip and equality at the joint types of pairs of generators,
transitivity at the iterated joint types of triples. Hence validity is
decidable.

\begin{theorem}[Enumeration up to isomorphism]
If the finite objects of a CEF $C$ can be listed up to isomorphism (finite
lists $L_0,L_1,\ldots$ such that every finite object has an isomorphic copy
in some $L_b$), then so can the finite objects of $\Stab(C)$: list, at
stage $B$, the valid tables with at most $B$ generators taken from
$L_0,\ldots,L_B$. Every code is isomorphic to a listed one: move its
generators to listed copies (rigidity makes the composites identities),
read off its table, and find an agreeing table in the list.
\end{theorem}

Starting from the finite ordinals this enumerates finite pre-dilators,
pre-ptykes of type~2, of type~3, \dots, up to isomorphism, and every element
of $\mathrm{FinDil}_C$ occurs in the lists. Since equality in
$\mathrm{FinDil}_C$ is decidable, the \emph{least} index hitting a given
element is a proposition and can be extracted from the truncation:

\begin{corollary}[Canonical representatives]
$\mathrm{FinDil}_C$ is equivalent to a decidable subset of $\mathbb{N}$
(the indices that are least in their class), and there is a section
$\mathrm{FinDil}_C\to\mathrm{Codes}$.
\end{corollary}

Restricted to monotone codes over $\LO$, whose monotonicity is checked at the
joint types of each generator with itself (decidably), this gives a
canonical sequence of codes of finite dilators. It replaces the parameter
of the $\Pi^1_4$ reduction: every code decoding to a dilator (it is then
monotone, Theorem~\ref{thm:mono}) is isomorphic to a member of the
sequence.

\paragraph{Uniqueness up to isomorphism.} In the proof of the main theorem,
the multicoproducts of two decoded codes must be unique \emph{strictly}:
two listed cospans through which a pair of maps factors must be equal
\emph{as codes}. This is achieved by a canonical form for the listed
cospans (which generators are kept, which comparisons are recorded, which
maps are the legs) and a determinacy argument showing that a factoring
cospan in canonical form is the covering one. The univalent statement is
uniqueness up to unique isomorphism, and it has a short generic proof, valid
in any category of embeddings. Call a cospan $u:X_1\to Z\leftarrow X_2:v$
\emph{jointly surjective} if every point of $Z$ comes from $X_1$ or from
$X_2$; listed sum codes have this property.

\begin{proposition}
Let $(u,v)$ be jointly surjective, and let a pair $(f,g)$ factor through it
by $m:Z\to Y$, and through any cospan $(u',v')$ by $m':Z'\to Y$. Then there
is a map $Z\to Z'$ compatible with the legs and with $m,m'$, and maps out
of $Z$ are determined by their composites with $u$ and $v$.
\end{proposition}

The map sends a point $u(a)$ to $u'(a)$ and $v(b)$ to $v'(b)$. It is well
defined because $m'$ is injective, and cartesian because $m$ is. With joint
surjectivity on both sides and rigid decoded objects, the two mediating
maps are inverse. Strict uniqueness is then recovered generically: if iso
of listed cospans is decidable (it is, the hom-sets being finite), drop
from the list every entry that has an isomorphic copy earlier in the list.
The new list still covers (the least isomorphic copy of a covering entry is
kept) and is strictly unique (two kept entries through which $(f,g)$
factors are isomorphic, so neither comes before the other). This gives a
second proof of the main theorem without the determinacy argument. The
canonical forms are still used, to know that the legs of the covering
cospan are the injections, but not their uniqueness.

\begin{remark}
In terms of size, the second proof is not shorter: the Stab-specific part
drops from about 290 to about 110 lines, but the generic parts (mediation
and deduplication) add about 330. The gain is conceptual. Uniqueness of
codes is replaced by universal properties, and, because finite objects are
rigid, uniqueness up to unique isomorphism becomes equality in the set
$\mathrm{FinDil}_C$.
\end{remark}

\section{What one may hope for}

\paragraph{A type-3 collapse.} Girard~\cite{GirH} conjectured that
$\Pi^1_n$-comprehension is equivalent to the totality of a canonical
``exponential'' $A_n$ of type $n+1$; known for $n=0$ ($2^X$) and $n=1$
($\Lambda$, or Freund's $\vartheta$~\cite{Fre19,AFW}). Freund showed that
fixed points of normal dilators correspond to $\Pi^1_1$-induction, and fixed
points of normal 2-ptykes to $\Pi^1_2$-induction~\cite{FrePilot}, while
Bachmann--Howard \emph{collapses} of dilators correspond to
$\Pi^1_1$-comprehension~\cite{Fre19}. The
natural candidate for $n=2$ is therefore the \emph{collapse of a 2-ptyx}:
for $P:\mathrm{DIL}\to\mathrm{DIL}$, a dilator $\Theta_P$ with a collapse
$\vartheta:P(\Theta_P)\Rightarrow\Theta_P$, so that $P$ is applied to the
dilator being built. Its terms are finite: $\vartheta(K,z,\vec y,\vec t)$
with $K$ a code, $z\in P(\mathrm{Decode}\,K)([c])$, $\vec y$ the arguments
and $\vec t$ the images of the generators of $K$ (its support). The
supports are \emph{functions} (trace terms with variables), and, following
Towsner~\cite{TowPoly}, two collapses are compared by a closure rule:
$\vartheta\sigma<\vartheta\tau$ if $\vartheta\sigma$ is not closed under
some support function $\gamma$ of $\tau$, i.e.\ $\gamma$ with its variables
placed just below $\vartheta\sigma$ is $\ge\vartheta\sigma$; symmetrically;
and otherwise $\sigma$ and $\tau$ are compared inside $P$ after selecting
the amalgam of the two codes from the multicoproduct list. Atoms lie below
all collapses; placing them above (as cardinals) is incompatible with the
closure rule for a total order. The
statement ``$\Theta_P$ is a dilator for every 2-ptyx $P$'' is $\Pi^1_4$,
the complexity of $\Pi^1_2$-$\mathrm{CA}_0$; no well-ordering principle for
$\Pi^1_2$-$\mathrm{CA}_0$ is known (compare Rathjen's conjecture on
$\beta$-models~\cite[Conj.~6.6]{Rat21}, a type-2 statement, and the ordinal
analysis~\cite{Rat05}).

What is established: a naive variant (collapses without free variables)
reduces to $\vartheta$ of an ordinary dilator and is provable in
$\Pi^1_1$-$\mathrm{CA}_0$; the candidate reduces to $\vartheta_D$ when $P$
lifts $D$; a single instance (iteration of $\Delta$ along $\omega$) should
exceed $\Pi^1_1$-$\mathrm{CA}_0$, modulo a comparison of Freund's
$\vartheta$-order with Buchholz's $\psi$-order~\cite{Buc86}.  Open: an informal proof
that $\Theta_P$ is well-founded for every 2-ptyx (some form of induction on
dilators), and the converse.

\paragraph{Towsner's polymorphic notations.} Towsner's ordinal
analyses~\cite{TowPMP,TowHalf,TowPoly} towards parameter-free
$\Pi^1_2$-comprehension use ordinal bounds with free
variables (functions that preserve well-foundedness, i.e.\ dilators) and
cardinals of function sort, whose coefficient sets contain functions that
are instantiated at the parameters of the other term. These are dilators
and type-2 ptykes in all but name; the semantics above could give them a
precise constructive meaning. His rule for function-valued coefficients
gives the closure rule used above; his placement of collapses below
variables, on the other hand, relies on his order being partial on open
terms and is not available for a dilator, whose order is total.

\paragraph{Other well-ordering principles.} $X\mapsto\varphi X0$
($\mathrm{ATR}_0$) and $X\mapsto\Gamma_X$ fit the same pattern as
$\varepsilon_X$.

\section{What has been implemented}

The development is written in a small dependently typed language and
checked by its type checker (about 26{,}800 declarations for the whole
library; the files below take a few seconds each). Main files, in the
\texttt{UNIT/} directory:
\begin{itemize}
\item categories of embeddings: \texttt{CatEmb}, \texttt{CatEmbFin} (CEF,
  Diers multicoproducts), \texttt{StabFun}, \texttt{StabPb},
  \texttt{StabDec} (decidable supports), \texttt{StabSupp} (least supports,
  unique normal forms), \texttt{StabCode}, \texttt{StabCover},
  \texttt{StabDir}, \texttt{StabFin}, \texttt{StabHom}, \texttt{StabSub},
  \texttt{StabAmal}, \texttt{StabRigid}, \texttt{StabJt0},
  \texttt{StabSum}, \texttt{StabJt}, \texttt{StabCEF} (main theorem);
\item base instances: \texttt{FinOrd}, \texttt{FinOrdEx},
  \texttt{LOCat}, \texttt{LOFin}, \texttt{LOCEF};
  well-foundedness hierarchy: \texttt{LOWF};
\item Girard's reductions: \texttt{KBTree}, \texttt{LabTree},
  \texttt{Girard}, \texttt{GirardPos}, \texttt{LabTree2}, \texttt{Girard2},
  and at type 3 \texttt{NSumLO}, \texttt{LabTree3}, \texttt{Girard3};
\item almost-full relations and the intuitionistic Ramsey theorem:
  \texttt{AlmostFull}; monotone finite pre-dilators are dilators:
  \texttt{FinDilWF};
\item examples: \texttt{DilLO}, \texttt{LOEx} (polynomials, $2^X$,
  towers, evaluations), \texttt{EpsTerm}, \texttt{EpsFun}, \texttt{EpsWF}
  ($\varepsilon_X$), \texttt{ThetaTerm}, \texttt{ThetaLO},
  \texttt{ThetaLO2}, \texttt{ThetaLO3}, \texttt{ThetaFun},
  \texttt{ThetaWF}, \texttt{ThetaFun2} ($\vartheta$ as a type-2 ptyx);
\item with univalence and replacement: \texttt{LOUnivalent},
  \texttt{FinLOSmall}, \texttt{SFUnivalent}, \texttt{SFLocSmall},
  \texttt{FinDil} (small set, decidable equality), \texttt{StabTab},
  \texttt{StabTabCov}, \texttt{StabEnum} (enumeration up to isomorphism),
  \texttt{MonoEnum} (canonical sequence of monotone codes),
  \texttt{FinDilCanon} (canonical representatives), \texttt{CospanUniq},
  \texttt{Dedup}, \texttt{StabCEF2} (second proof of the main theorem),
  \texttt{DilSet} (dilators form a set), \texttt{Gaunt}, \texttt{DilCat}
  (monotone dilators as a category of embeddings; type-2 ptykes form a set),
  \texttt{UnionWF}; classical completion of the $\Pi^1_2$ theorem:
  \texttt{GirardClassical}; dilators are monotone: \texttt{DilMono};
\item type 3 (in progress): \texttt{ThetaP3Term} (terms of $\Theta_P$ and a
  decidable comparison, defined with a fuel bound); a Python mirror of the candidate in
  \texttt{mirror/}.
\end{itemize}
An earlier, more hands-on development over finite ordinals
(\texttt{CatDil}, \texttt{DilPtyx}, \texttt{Ptyx*}, \texttt{Ptyx3*}) treats
ptykes of types 2 and 3 directly, with germs, subptykes and directed
colimits; it is subsumed by the abstract theorem.

\begin{thebibliography}{99}
\bibitem{AR} B.~Afshari and M.~Rathjen, Reverse mathematics and
  well-ordering principles: a pilot study, \emph{Ann.\ Pure Appl.\ Logic}
  160 (2009) 231--237.
\bibitem{AFW} J.~P.~Aguilera, A.~Freund and A.~Weiermann, Induction on
  dilators and Bachmann--Howard fixed points, arXiv:2412.13051.
\bibitem{Buc86} W.~Buchholz, A new system of proof-theoretic ordinal
  functions, \emph{Ann.\ Pure Appl.\ Logic} 32 (1986) 195--207.
\bibitem{CoqRamsey} Th.~Coquand, A direct proof of Ramsey's Theorem, note,
  2010, \url{https://www.cse.chalmers.se/~coquand/ramsey1.pdf}.
\bibitem{Coq89} Th.~Coquand, Categories of embeddings, \emph{Theoret.\
  Comput.\ Sci.} 68 (1989) 221--237.
\bibitem{Die80} Y.~Diers, Cat\'egories localement multipr\'esentables,
  \emph{Arch.\ Math.} 34 (1980) 344--356.
\bibitem{Fre19} A.~Freund, $\Pi^1_1$-comprehension as a well-ordering
  principle, \emph{Adv.\ Math.} 355 (2019) 106767.
\bibitem{FrePilot} A.~Freund, Well ordering principles and
  $\Pi^1_4$-statements: a pilot study, arXiv:2006.12111.
\bibitem{Gir81} J.-Y.~Girard, $\Pi^1_2$-logic, Part 1: Dilators,
  \emph{Ann.\ Math.\ Logic} 21 (1981) 75--219.
\bibitem{GirH} J.-Y.~Girard, Herbrand's theorem and proof theory, in:
  \emph{Proceedings of the Herbrand Symposium (Marseille 1981)},
  North-Holland, 1982, 29--38.
\bibitem{GirPTLC} J.-Y.~Girard, \emph{Proof Theory and Logical
  Complexity}, vol.~1, Bibliopolis, 1987; vol.~2, manuscript.
\bibitem{GN} J.-Y.~Girard and D.~Normann, Embeddability of ptykes,
  \emph{J.\ Symbolic Logic} 57 (1992) 659--676.
\bibitem{MM} A.~Marcone and A.~Montalb\'an, The Veblen functions for
  computability theorists, \emph{J.\ Symbolic Logic} 76 (2011) 575--602.
\bibitem{Osm1} A.~Osmond, On Diers's theory of spectrum I: stable functors
  and right multi-adjoints, arXiv:2012.00853.
\bibitem{Osm2} A.~Osmond, On Diers's theory of spectrum II: geometries and
  dualities, arXiv:2012.02167.
\bibitem{Rat05} M.~Rathjen, An ordinal analysis of parameter free
  $\Pi^1_2$-comprehension, \emph{Arch.\ Math.\ Logic} 44 (2005) 263--362.
\bibitem{Rat21} M.~Rathjen, Well-ordering principles in proof theory and
  reverse mathematics, arXiv:2010.12453.
\bibitem{TowPMP} H.~Towsner, Proofs that modify proofs, arXiv:2403.17922.
\bibitem{TowHalf} H.~Towsner, Proofs that modify proofs, 1/2,
  arXiv:2506.16491.
\bibitem{TowPoly} H.~Towsner, Polymorphic ordinal notations,
  arXiv:2504.02131.
\bibitem{HoTT} The Univalent Foundations Program, \emph{Homotopy Type
  Theory: Univalent Foundations of Mathematics}, Institute for Advanced
  Study, 2013.
\bibitem{Rij} E.~Rijke, The join construction, arXiv:1701.07538.
\bibitem{Win} G.~Winskel, Event structures, Cambridge University Computer
  Laboratory, Technical Report No.~95.
\end{thebibliography}

\end{document}
